\documentclass[10pt,a4paper]{amsart}
\usepackage[margin=19mm]{geometry}
\usepackage{amsmath,amssymb,mathtools}
\usepackage[scr=rsfs,cal=euler]{mathalfa}
\usepackage{xfrac}
\usepackage[unicode,hidelinks,bookmarks=false]{hyperref}
\newcommand{\ie}{i.e.\ }
\newcommand{\cf}{cf.\ }
\newcommand{\resp}{respectively\ }
\newcommand{\Subsection}{\subsection}
\providecommand{\nobreakdash}{\nobreak\hbox{-}}
\setlength{\emergencystretch}{2em}
\multlinegap0pt
\usepackage[shortlabels]{enumitem}
\usepackage{amssymb}
\usepackage{mathtools}
\usepackage{float}
\newtheorem{lemma}{Lemma}
\newcommand{\amp}{\mathop{\:\:\,}\nolimits}
\newcommand{\prox}{\operatorname{prox}}
\title{Operator Splitting with Hamilton-Jacobi-based Proximals}
\author{Nicholas Di, Eric C. Chi, Samy Wu Fung}
\date{Source: arXiv:2601.22370v4 (CC BY 4.0)}
\begin{document}
\maketitle

\noindent\emph{Source and licence.} Operator Splitting with Hamilton-Jacobi-based Proximals. Version arXiv:2601.22370v4 (CC BY 4.0), distributed by arXiv under the Creative Commons Attribution 4.0 International licence, http://creativecommons.org/licenses/by/4.0/, as recorded in the arXiv licence metadata for this exact version and shown on the abstract page https://arxiv.org/abs/2601.22370v4. Full licence notice: \texttt{licenses/arxiv-2601.22370-CC-BY-4.0.txt}.\par\smallskip

\noindent\textit{Editorial scope.} Counted unit: the article's lemma PGD\_Convergence (source0.tex lines 1064--1070). The article's complete original proof of it is delivered with it (source0.tex lines 1072--1111, one \texttt{proof} environment of 40 lines). No mathematics is written, repaired or re-derived: the statement and the proof are the article's own lines, in the article's own notation, and no retained line is substituted or re-flowed.  No further source-local context is needed: every cross-reference of every retained line resolves inside the two delivered spans.  Nothing was omitted from any retained span, so the package carries no omission record.  No retained line contains a citation, so the wrapper carries no bibliography and no bibliography entry.  No retained line contains a figure environment, an image-inclusion command or drawing code, so no image is delivered and the package declares no asset dependency.  Editorial formatting only, in addition: an AMS \texttt{amsart} preamble that restates the article's own declarations and macros used by the retained lines, namely the environments \texttt{lemma} on separate counters, together with the standard packages those lines need.  The retained environments print in this document as Lemma 1 in order of appearance, whereas the article prints the same items with the numbering of its own full document.  The complete native source of the article as distributed in the arXiv e-print for this exact version is bundled unchanged as \texttt{sources/193651/source0.tex}.
\begin{lemma}[PGD Diminishing Stepsizes]\label{lem: PGD_Convergence}
Consider $F = f+ g$ where $f$ is convex with $L$-Lipschitz continuous gradient, $g$ is proper, convex and lsc, and the solution set $S = \arg\min F$ is nonempty. The proximal gradient iterations
\begin{eqnarray}
    x_{k+1} &=& \prox_{t_k g}(x_k - t_k \nabla f(x_k))
\end{eqnarray}
with $t_k \leq 1/L$ and $\sum_{k=0}^{\infty} t_k = \infty$ then every cluster point of $(x_k)_{k \geq 0}$ lies in $S$.
\end{lemma}

\begin{proof}
Let $u\in\mathbb{R}^n$. Using the proximal optimality condition for $g$ at
$x_{k+1}=\prox_{t_k g}(x_k-t_k\nabla f(x_k))$, convexity of $f$, and the $L$-smoothness of $f$,
we obtain the standard proximal-gradient three-point inequality
\begin{eqnarray}\label{eq:PGD_TPI}
    2t_k(F(x_{k+1}) - F(u))
    &\leq & \|x_k-u\|^2 - \|x_{k+1}-u\|^2 - (1-Lt_k)\|x_{k+1}-x_k\|^2 .
\end{eqnarray}
Taking $u \in S$ (so $F(u)=F^*$) and using $t_k\le 1/L$ gives
\begin{eqnarray}
2t_k(F(x_{k+1})-F^*) &\le& \|x_k-u\|^2-\|x_{k+1}-u\|^2,
\end{eqnarray}
hence $\|x_{k+1}-u\|\le \|x_k-u\|$. Thus $(x_k)$ is Fej\'er monotone with respect to $S$,
hence bounded with at least one cluster point. Summing \eqref{eq:PGD_TPI} over $k$ and dropping the
nonnegative term $(1-Lt_k)\|x_{k+1}-x_k\|^2$ yields
\begin{eqnarray}\label{eq:summable_pgd}
    \sum_{k=0}^{\infty} t_k(F(x_{k+1}) - F^*) &\leq& \frac{1}{2}\|x_0-u\|^2 < \infty.
\end{eqnarray}

We take $u=x_k$ in \eqref{eq:PGD_TPI} 
\begin{eqnarray}
    2t_k(F(x_{k+1})-F(x_k)) &\le& -(1-Lt_k)\|x_{k+1}-x_k\|^2 \le 0,
\end{eqnarray}

so $F(x_k)$ is non-increasing, hence converges to some $L \geq F^*$. If $L > F^*$, then
$F(x_{k+1}) - F^* \geq L - F^* > 0$ for all sufficiently large $k$, which combined with
$\sum_{k=0}^\infty t_k=\infty$ implies
\begin{eqnarray}
\sum_{k=0}^{\infty} t_k(F(x_{k+1}) - F^*) &\geq& (L-F^*)\sum_{k=0}^{\infty} t_k \amp=\amp \infty,
\end{eqnarray}
contradicting \eqref{eq:summable_pgd}. Thus $\lim_{k \to \infty} F(x_k) = F^*$.

Finally, let $\bar{x}$ be any cluster point with $x_{k_n} \to \bar{x}$. Since $f$ is continuous and
$g$ is lsc, $F=f+g$ is lsc, hence
\begin{eqnarray}
    F(\bar{x}) &\leq& \liminf_{n \to \infty} F(x_{k_n}) \amp=\amp F^*,
\end{eqnarray}

which implies $\bar{x} \in S$.
\end{proof}

\end{document}
